\documentclass[11pt,a4paper]{article}
\usepackage{amsmath,amssymb,amsthm}
\usepackage{hyperref}
\usepackage{graphicx}

\title{On the Terminal Spectrum of Weierstrass Prym Eigenforms\\
in the Minimal Stratum of Genus~4}
\author{}
\date{}

\newtheorem{theorem}{Theorem}
\newtheorem{proposition}[theorem]{Proposition}
\newtheorem{corollary}[theorem]{Corollary}
\newtheorem{remark}[theorem]{Remark}

\begin{document}
\maketitle

\begin{abstract}
Weierstrass Prym eigenforms \(W_D(2g-2)\) exist only for \(g\le 4\).
In genus~4 the entire family \(W_D(6)\) is contained in the even component
of the minimal stratum \(\Omega\mathcal{M}_4(6)\). Combined with the
Weierstrass filtration of Yu--Zuo, this forces the individual Lyapunov
exponents on the positive Prym block \(H_1^+\) to be
\[
\lambda_1^+=\tfrac47,\qquad\lambda_2^+=\tfrac27
\]
for every admissible discriminant \(D\). The odd component of the same
stratum carries the contrasting rigid spectrum \(3/7,2/7,1/7\). The result
is an intersection of existing theorems; no new computation is claimed.
\end{abstract}

\section{Finiteness of the family}
Weierstrass Prym eigenforms in the minimal stratum exist only in low genus.
\begin{proposition}[Guti\'errez-Romo--Pardo et al.]
The loci \(\Omega E_D(2g-2)\) are empty for all \(g\ge 5\).
\end{proposition}
Consequently genus~4 is the terminal case in which these dynamics appear.

\section{Spin parity and structural containment}
\begin{proposition}[M\"oller]
Let \((X,\omega)\in W_D(6)\) with \(\operatorname{div}(\omega)=6P\).
The spin structure determined by \(\mathcal{O}_X(3P)\) is even.
Equivalently,
\[
W_D(6)\;\subset\;\Omega\mathcal{M}_4(6)^{\mathrm{even}}.
\]
\end{proposition}

\section{Weierstrass filtration on the minimal stratum}
\begin{theorem}[Yu--Zuo]
For every Teichm\"uller curve in a connected component of
\(\Omega\mathcal{M}_g(2g-2)\) with \(g\le 5\), the Weierstrass gap sequence
is constant and the Hodge bundle splits into a direct sum of line bundles
whose degrees yield the following exponents \(w_i=k_i/(2g-1)\):

\begin{center}
\begin{tabular}{lcc}
Component & Gap sequence \(G_p\) (\(g=4\)) & Exponents \\
\hline
even & \(\{1,2,4,7\}\) & \(1,\ 4/7,\ 2/7,\ 1/7\) \\
odd  & \(\{1,2,3,7\}\) & \(1,\ 3/7,\ 2/7,\ 1/7\) \\
hyperelliptic & \(\{1,3,5,7\}\) & \(1,\ 5/7,\ 3/7,\ 1/7\)
\end{tabular}
\end{center}
\end{theorem}

\section{Terminal spectrum}
\begin{corollary}
For every admissible discriminant \(D\), every Teichm\"uller curve in the
Weierstrass Prym locus \(\Omega E_D(6)\) satisfies
\[
\lambda_1^+=\frac47,\qquad\lambda_2^+=\frac27
\]
on the \(4\)-dimensional block \(H_1^+\). (The complementary block \(H^-\)
carries the known exponents \(\pm 1/7\).)
\end{corollary}

The sum \(\lambda_1^++\lambda_2^+=6/7\) is the EKZ formula; the individual
split is forced by the even-component filtration of Yu--Zuo together with
M\"oller's parity containment.

\begin{remark}[Contrast]
The odd component of \(\Omega\mathcal{M}_4(6)\) is likewise rigid, but with
the distinct spectrum \(3/7,2/7,1/7\). The Prym--Weierstrass family never
enters that component.
\end{remark}

\begin{figure}[ht]
\centering
\includegraphics[width=0.85\textwidth]{spectral_contrast_g4.pdf}
\caption{Weierstrass exponents on the three components of
$\Omega\mathcal{M}_4(6)$ (Yu--Zuo). All $W_D(6)$ sit in the even
component and therefore inherit $4/7,2/7$.}
\end{figure}


\section{Attribution}
\begin{itemize}
\item Existence ceiling \(g\ge 5\): Guti\'errez-Romo--Pardo et al.
\item Even spin of \(W_D(6)\): M\"oller.
\item Rigid Weierstrass exponents on even/odd components: Yu--Zuo
  (arXiv:1203.6053).
\item Sum formulae and \(H^-=\pm 1/7\): Eskin--Kontsevich--Zorich / M\"oller.
\end{itemize}
No computational determination of the individual exponents is claimed.

\end{document}
